\subsection{Eleven-element corners and the seventeen-arrow act gates}

\begin{lemma}[Proved: two-dimensional local module pairs]
\label{CP2:A-local-two-by-two}
Let $A$ be a finite-dimensional local rational algebra and let
$W,V$ be right and left modules of dimension two. They cannot have
\[
 (\dim\operatorname{Tor}_0^A(W,V),\ldots,
   \dim\operatorname{Tor}_3^A(W,V))=(1,0,1,0).
\]
Nor can they have bounded Tor with tensor dimension two,
zero first Tor, and nonzero second Tor.
\end{lemma}
\begin{proof}
If either module is semisimple, the local minimal-resolution
argument in Lemma~\ref{A:local-target-facts} makes its positive
Tor with a nonprojective opposite module nonzero in every
positive degree. A projective opposite module gives no positive
Tor. Thus both modules are nonsemisimple. Their residue division
algebra is then $\mathbb Q$, and both are cyclic. A cyclic
projective cover of $V$ has kernel $I$, of dimension $\dim A-2$.
We may suppose $I\ne0$, since otherwise $V$ is projective.
The first Tor vanishing gives
\[
 0\longrightarrow W\otimes_A I\longrightarrow W
   \longrightarrow W\otimes_A V\longrightarrow0.
\]
The tensor product of two nonzero modules over a finite local
algebra is nonzero: it surjects onto the tensor product of their
tops over the residue division algebra. This already excludes
tensor dimension two.

For tensor dimension one, $W\otimes I$ has dimension one, so
$I$ has one-dimensional top and is cyclic. Its cyclic cover
has kernel $L=\Omega^2 V$ of dimension two. If $L$ is semisimple,
then $\operatorname{Tor}_3(W,V)=\operatorname{Tor}_1(W,L)$
is nonzero, since $W$ is nonprojective. Otherwise $L$ is cyclic.
The first three projectives in a minimal resolution of $V$
are consequently $A$. After tensoring with $W$, the first two
differentials have ranks one and one. The required second Tor
dimension one forces the third differential to be zero.
The next projective is $A^t$ with $t>0$: otherwise $L$ is
projective, contrary to its dimension two and $\dim A>2$.
The fourth differential has image in $(WJ)^t$, a proper
subspace of $W^t$, by minimality. It cannot cover the whole
kernel of the zero third differential. Hence third Tor is
nonzero, a contradiction.
\end{proof}

\begin{lemma}[Proved: the five-dimensional two-by-two alternatives]
\label{CP2:A-five-algebra-two-by-two}
Let $A$ be a nonlocal rational algebra of dimension five,
$\dim W=\dim V=2$, and suppose its Tor is bounded,
$\operatorname{Tor}_1(W,V)=0$, and
$\dim\operatorname{Tor}_2(W,V)=1$.
If the tensor dimension is one, there is no such pair.
If it is two, the only remaining case in the classification
of Lemma~\ref{A:nonlocal-five-target} is the zero-composite
three-vertex line algebra, with
\[
 V=S_1\oplus S_3,\qquad W=T_1\oplus T_3.
\]
The two-dimensional space of balanced forms for this pair
has a basis consisting of matrices of rank one.
\end{lemma}
\begin{proof}
The persistence alternatives in that lemma contradict boundedness.
For $A=B\times\mathbb Q$, with $\dim B=4$ local, all positive
Tor is on the $B$ components. A one-dimensional nonzero
component forces persistence by the local minimal-resolution
argument. Thus both $B$ components have dimension two,
leaving no other components. Tensor dimension one is excluded
by Lemma~\ref{A:local-four-small-target}, and tensor dimension
two by Lemma~\ref{CP2:A-local-two-by-two}.

For the line algebra, use the five left and five right
indecomposables and their displayed resolutions in
Proposition~\ref{A:three-by-six-CP2}. Nonzero second Tor
requires a $T_1$ summand on the right and an $S_3$ summand
on the left. Each must occur exactly once. Neither remaining
one-dimensional summand can be the middle-vertex simple,
since that would give nonzero first Tor; repeating $T_1$
or $S_3$ would give second Tor dimension two. Thus the
pair is exactly the displayed pair. Its tensor dimension
is two, with one coordinate each at vertices1 and3.
The corresponding two coordinate functionals are rank-one
matrices and span the balanced-form space.

For the two-vertex exceptional algebra, the two-dimensional
left modules are
$S_1^{\oplus2},S_1\oplus S_2,S_2^{\oplus2},L,Ae_2$.
The first and last give no second Tor. For $S_2^{\oplus2}$
the second Tor dimension is even. For $S_1\oplus S_2$,
every right summand contributing to second Tor also gives
nonzero first Tor. With $V=L$, the only permissible
contributor is one copy of $T_2$: $T_1$ and $R$ give
nonzero first Tor. A right module of dimension two cannot
contain just one $T_2$ and a projective complement, since
the two indecomposable projectives have dimensions three
and two. Two copies of $T_2$ give second Tor dimension
two. This proves the exclusion. All these assertions also
follow by summing the independently checked literal
$5\times5$ Tor tables in
\path{monoid/certificates/five_dimensional_pairs.json}.
\end{proof}

\begin{proposition}[Proved: rational three-state cones through order eleven]
\label{CP2:A-eleven-three-state-rational}
An ordinary act cone over a monoid of order at most eleven,
with at most three states, cannot have the rational homology
and cohomology ring of $\CP^2$.
\end{proposition}
\begin{proof}
Orders through ten follow from
Proposition~\ref{CP2:A-ten-six-state-rational}.
Assume the corner order is eleven. The relative-cup and
two-row obstructions require at least three minimal rows
and two minimal columns. A nontrivial minimal group would
cost at least thirteen elements, so the ideal is a band.
If its action on the at most three states is not constant,
Lemma~\ref{CP2:A-constant-quotient-H2} pulls every second
class from a cone with at most two states, where every
second product is zero. We may therefore assume constant
band action. Fewer than three states are already excluded.

Compression gives a monoid with minimal band $I\times S$
and the same outside algebra $A$ as the original corner.
For at most two outside elements use
Lemma~\ref{CP2:A-two-outside-acts}. The only remaining
original band shapes and outside counts are
$3\times2$ with five outside elements and
$4\times2$ with three. In the second case the compressed
row and column augmentation dimensions are three and two;
Lemma~\ref{CP2:A-local-three-small-module} and
Lemma~\ref{A:nonlocal-four} exclude bounded nonzero
second Tor. In the first case $\dim A=5$ and the two
augmentation modules both have dimension two. The required
Tor dimensions are $(1,0,1,0,\ldots)$. The preceding two
lemmas exclude these dimensions for local and nonlocal
$A$, respectively. All comparisons are actual maps or
homotopy equivalences, so their cohomology pullbacks
preserve products.
\end{proof}

\begin{theorem}[Proved: eleven-element corners with five states]
\label{CP2:A-eleven-five-state-actual}
An ordinary act cone over a monoid of order at most eleven,
with at most five states, cannot have type $\CP^2$.
\end{theorem}
\begin{proof}
Orders through ten are excluded by the rational theorem;
through three states use the preceding proposition, and
four states use Corollary~\ref{CP2:A-four-state-small}.
Thus $|M|=11$ and $|S|=5$. The same row, column and
group budgets reduce to a rectangular minimal band.
A band-fixed state is globally fixed and kills second
products. Without such a state, the only shapes not
already excluded by at most two outside elements are
$3\times2$ with five outside elements and
$4\times2$ with three.

Let $o$ be the number of two-point core orbits. At five
states $o=1$ or2. For $o=1$, every band actor is constant:
its value on the sole core orbit is constant, and
idempotence determines its value on every transient state.
Compression therefore gives an actual monoid model of
order20 in the $3\times2$ case, contradicting the monoid
minimum. In the $4\times2$ case its order is23 and its
minimal band is $4\times5$, contrary to the complete
classification at order23.

For $o=2$, Lemma~\ref{CP2:A-two-column-orbit-kernel}
gives a compressed monoid $N$ with $H_2(BN;\mathbb Q)
=\mathbb Q^2$, $H_4(BN;\mathbb Q)=\mathbb Q$, no other
positive homology, and degree-two cup form of rank one.
Its column augmentation module has dimension two.
For the $4\times2$ shape, the outside algebra has
dimension three and the row module has dimension three;
the small-local and nonlocal-four lemmas again force
second Tor to vanish.

For the $3\times2$ shape, the outside algebra has
dimension five and both modules have dimension two.
Their Tor dimensions are $(2,0,1,0,\ldots)$.
The local case is excluded by
Lemma~\ref{CP2:A-local-two-by-two}; the nonlocal case
leaves only the line-algebra pair in
Lemma~\ref{CP2:A-five-algebra-two-by-two}.
Each of its two basis balanced forms has matrix rank one
and hence square zero by Lemma~\ref{A:rank-one-square}.
The symmetric cup matrix in this basis is therefore
$\left(\begin{smallmatrix}0&c\\c&0\end{smallmatrix}\right)$.
Over $\mathbb Q$ its rank is zero or two, never one.
This contradicts the quotient comparison's rank-one cup
form, and finishes every case.
\end{proof}

\begin{proposition}[Proved: the eleven-by-two four-arrow collage]
\label{CP2:A-eleven-two-four}
A two-object category with corner orders eleven and two
and four connecting arrows cannot have type $\CP^2$.
\end{proposition}
\begin{proof}
By opposition let the two-element corner act on the left
of a right-$M$ act $S$, where $|M|=11$ and $|S|=4$.
The two-element monoids are the group $C_2$ and
$\{1,e\}$ with $e^2=e$.
For the latter, put $T=eS$. Sending $S$ to $T$ and
both corner elements to the identity defines a functor
to the ordinary $M$-act cone on $T$. Its inclusion is a
section. The components $1_M$ and $e$ give a natural
transformation from the identity of the original category
to the inclusion followed by the quotient: on a cross
arrow the equation is precisely $es=es$.
Thus this is a homotopy equivalence, excluded by the
three-state rational or four-state actual theorem.

For the group $C_2$, exact rational orbit comparison
gives a rational homology and cohomology-ring isomorphism
to the ordinary $M$-act cone on $S/C_2$. If the action
is nontrivial, this orbit set has at most three elements,
contradicting Proposition~\ref{CP2:A-eleven-three-state-rational}.
If the action is trivial, the ordinary cone on all four
states is an actual categorical retract: quotient the
group corner to the identity, and include that identity
as a section. The quotient is still a rational homology
isomorphism. A retract of a simply connected space is
simply connected, and its integral homology groups are
direct summands of those of $\CP^2$. The rational
isomorphism consequently forces integral homology
isomorphisms in every degree. Naturality gives the
primitive square, so this retract has type $\CP^2$ by the same integral recognition argument as in
Theorem~\ref{A:monoid-upper23}. This contradicts
the four-state act theorem. Both monoids are excluded.
\end{proof}


\begin{corollary}[Proved: all act cones through seventeen arrows]
\label{CP2:A-act-cones-through-seventeen}
No ordinary finite-monoid act cone with at most seventeen
arrows has type $\CP^2$.
Every two-object category of at most seventeen arrows with one
corner of order two is excluded as well.
\end{corollary}
\begin{proof}
For an ordinary cone write $m$ for the corner order and $s$
for the positive number of states, so $m+s+1\le17$.
Orders through nine use the bounded-nine theorem; order ten
has $s\le6$ and uses the ten-state-budget rational theorem;
order eleven has $s\le5$ and uses
Theorem~\ref{CP2:A-eleven-five-state-actual}. For order twelve,
$s\le4$, use the four-state theorem or the three-state minimum.
For $m\ge13$, there are at most three states, so the universal
one- and two-state obstructions and the exact three-state
minimum twenty-two apply.

A disconnected category is excluded by connectedness. If both cross directions are nonempty, Proposition~\ref{CP2:A-baseline-strong-two} applies. In the remaining triangular case, take opposition if necessary. For a corner of order two, $m+s+2\le17$. One, two and three
connecting arrows are covered by the same universal cross-arrow
results. If $s\ge4$, then $m\le11$. Exact rational orbit
comparison at the order-two corner preserves the rational
homology and ring, with at most $s$ states. The bounded-nine
and ten-six-state theorems cover $m\le10$. The only remaining
budget is $(m,2,s)=(11,2,4)$, which is the preceding proposition.
\end{proof}
